\section{The GeoSection Algorithm}
\label{sec:algorithm}

\subsection{Setup and Notation}

Let $G = (V, E, w)$ be the state's census-tract adjacency graph.
Vertices $V$ are census tracts; edges $E$ connect tracts sharing a common
boundary; edge weights $w(e) > 0$ are shared boundary lengths in metres,
extracted from TIGER geometries.
Vertex weights $\mathbf{p} \in \mathbb{Z}_{>0}^{|V|}$ are census tract populations.
Let $k$ be the number of congressional districts allocated to this state.

\begin{definition}[Feasible bisection]
A \emph{feasible bisection} of $G$ with ratio $i:(k-i)$, $1 \leq i \leq \lfloor k/2 \rfloor$,
is a partition $(S, V \setminus S)$ satisfying the population balance constraint:
\[
  \frac{\sum_{v \in S} p_v}{\sum_{v \in V} p_v} \in
  \left[\frac{i/k}{1 + \delta},\; \frac{i/k \cdot (1+\delta)}{1}\right]
\]
for balance tolerance $\delta$ (congressional default: $\delta = 0.005$).
The \emph{edge-cut} of a bisection is $\mathrm{EC}(S) = \sum_{e \in \delta(S)} w(e)$
where $\delta(S)$ is the set of edges crossing the boundary.
\end{definition}

\begin{definition}[Natural ratio]
The \emph{natural ratio} of $G$ for $k$ districts is the ratio $i^*:(k-i^*)$ where
\[
  i^* = \argmin_{1 \leq i \leq \lfloor k/2 \rfloor}
  \frac{\mathrm{EC}_{\min}(i:k-i)}{\sqrt{\min(i,\,k-i)}}
\]
and $\mathrm{EC}_{\min}(i:k-i) = \min_{\text{feasible } S} \mathrm{EC}(S)$
is the minimum edge-cut over all feasible bisections at ratio $i:(k-i)$.
The denominator $\sqrt{\min(i,k-i)}$ is the \emph{isoperimetric correction}.
\end{definition}

\subsection{Isoperimetric Normalisation}
\label{sec:normalisation}

The isoperimetric correction is the key innovation of GeoSection.
We motivate it formally.

\begin{remark}[Motivating heuristic: isoperimetric scaling]
\label{lem:isoperimetric}
For a disk-shaped uniformly populated region of total boundary length $L$,
the minimum-boundary-length cut separating fraction $f \in (0,1)$ of the
population satisfies approximately:
\[
  \mathrm{EC}_{\min}(f) \approx L \cdot \sqrt{\min(f, 1-f)}.
\]
This is a \emph{motivating heuristic}, not a theorem.
For $f$ small, enclosing fraction $f$ of a disk's area requires a circle of
perimeter $\propto \sqrt{f}$, giving the $\sqrt{f}$ scaling.
For $f$ near $1/2$, the optimal cut for a disk approaches a chord of length
roughly proportional to $\sqrt{1-f}$, consistent with the same formula.
For general convex regions --- especially non-circular shapes like US states ---
the scaling is approximate.
For elongated states (e.g., California), the cut length near $f = 1/2$ scales
linearly with the minor axis width, not as $\sqrt{f}$.
The normalisation $\mathrm{EC}/\sqrt{\min(i,k-i)}$ remains a principled
correction even when the approximation is not tight: it reduces the structural
advantage that small-$f$ splits enjoy over large-$f$ splits purely because
their boundaries are intrinsically shorter.
We use METIS-reported actual edge-cuts (not approximated minimums), so the
normalised comparison is between real boundary lengths, with the heuristic
correction applied uniformly.
\end{remark}

\begin{remark}
The isoperimetric correction $\sqrt{\min(i,k-i)}$ is a principled but approximate
normalisation (see Remark~\ref{lem:isoperimetric}).
Real states have highly non-uniform population density.
The correction remains meaningful in the non-uniform case:
$\sqrt{\min(i,k-i)}$ reduces the structural advantage of small-$i$ splits by
partially correcting for the fact that smaller-fraction boundaries are intrinsically
shorter regardless of compactness.
The correction's purpose is to prevent $1:(k-1)$ from winning by default;
it does not guarantee that the winner is the globally optimal ratio in
any formal sense.
\end{remark}

Without normalisation, the caterpillar pathology arises because:
\[
  \mathrm{EC}(1:k-1) \ll \mathrm{EC}(k/2:k/2)
\]
for any state with a compact urban core.
With normalisation:
\[
  \frac{\mathrm{EC}(1:k-1)}{\sqrt{\min(1,k-1)}} = \mathrm{EC}(1:k-1)
  \quad\text{vs.}\quad
  \frac{\mathrm{EC}(k/2:k/2)}{\sqrt{k/2}}
\]
The $1:k-1$ ratio loses its structural advantage.
For $k=8$ (Wisconsin), the normalised comparison is
$\mathrm{EC}(1:7)$ vs.\ $\mathrm{EC}(4:4)/\sqrt{4} = \mathrm{EC}(4:4)/2$.
If $\mathrm{EC}(1:7) < \mathrm{EC}(4:4)/2$, Milwaukee is genuinely compact
by the isoperimetric standard and the $1:7$ split wins.
If not, a more equal split wins.

\subsection{The GeoSection Algorithm}
\label{sec:geosection-alg}

\begin{algorithm}[H]
\caption{GeoSection($G$, $k$, $N$, $\delta$, $\lambda$)}
\label{alg:geosection}
\begin{algorithmic}[1]
\Require Graph $G = (V, E, w)$, districts $k \geq 1$, seeds per ratio $N$,
  balance tolerance $\delta$, directional penalty $\lambda \geq 0$
\Ensure Partition $\{D_1, \ldots, D_k\}$ of $V$

\If{$k = 1$} \Return $\{V\}$ \EndIf

\State $\mathrm{node\_ufactor} \gets 1 + \delta / k$
\hfill\Comment{per-level tolerance}

\State\Comment{\textbf{Phase 1: Ratio scan (isoperimetric normalisation)}}
\State $\mathrm{best\_norm} \gets +\infty$;
   $i^* \gets 1$; $S^* \gets \emptyset$

\For{$i = 1$ to $\lfloor k/2 \rfloor$}
  \State $\mathrm{EC}_{\min} \gets +\infty$; $S_{\min} \gets \emptyset$
  \For{$s = 1$ to $N$}
    \State $S \gets \mathrm{METIS}(G,\, i:(k-i),\, \mathrm{ufactor}=\mathrm{node\_ufactor},\, \mathrm{seed}=s)$
    \If{$\mathrm{EC}(S) < \mathrm{EC}_{\min}$}
      $\mathrm{EC}_{\min} \gets \mathrm{EC}(S)$; $S_{\min} \gets S$
    \EndIf
  \EndFor
  \State $\mathrm{norm}_i \gets \mathrm{EC}_{\min} / \sqrt{\min(i, k-i)}$
  \If{$\mathrm{norm}_i < \mathrm{best\_norm}$}
    $\mathrm{best\_norm} \gets \mathrm{norm}_i$; $i^* \gets i$; $S^* \gets S_{\min}$
  \EndIf
\EndFor

\State\Comment{\textbf{Phase 2: Directional penalty (if $\lambda > 0$)}}
\If{$\lambda > 0$}
  \State $\hat{u} \gets \mathrm{PCA\_minor\_axis}(\mathrm{centroids}(V))$
  \hfill\Comment{minor = narrowest direction}
  \State $w'(u,v) \gets w(u,v) \cdot (1 + \lambda \cdot |\!\sin(\theta(u,v))|)$
  \hfill\Comment{$\theta$ = angle between edge and cut direction}
  \State Re-run ratio scan on $G' = (V, E, w')$ to select $S^*$
\EndIf

\State\Comment{\textbf{Recursion}}
\State $V_L \gets S^*$;\ $V_R \gets V \setminus S^*$
\State $G_L \gets G[V_L]$;\ $G_R \gets G[V_R]$
\hfill\Comment{induced subgraphs}

\State \Return $\mathrm{GeoSection}(G_L, i^*, N, \delta, \lambda)$
         $\cup$ $\mathrm{GeoSection}(G_R, k - i^*, N, \delta, \lambda)$
\end{algorithmic}
\end{algorithm}

\paragraph{Key properties.}
\begin{itemize}
\item Each recursive call finds its own natural ratio independently.
  The right half's natural ratio is not constrained to equal the left half's.
\item The directional penalty is computed from the \emph{current subregion's}
  centroids, not the original state.
  After the first split, the left subregion's minor axis may differ from the
  full state's minor axis, allowing finer-grained orientation tracking.
\item For $k=2$, only one ratio ($1:1$) is possible; the algorithm reduces
  to standard minimum-edge-cut bisection.
\item For $k=1$, no METIS call is needed; the trivial single-district
  partition is returned.
\end{itemize}

\subsection{Directional Penalty}
\label{sec:phase2}

The directional penalty $\lambda$ in Phase 2 produces straighter district
boundaries by discouraging cuts that run \emph{parallel} to the subregion's
minor axis.
The motivation is geometric: a cut parallel to the minor axis produces a
narrow strip (like a salamander's tail); a cut perpendicular to the minor
axis bisects the subregion across its widest dimension.

\begin{definition}[Directional penalty weight]
Let $\hat{u}$ be the minor-axis direction of the current subregion
(unit vector in $\mathbb{R}^2$, computed by PCA on tract centroids).
The directional penalty re-weights each edge $(u,v)$ by:
\[
  w'(u,v) = w(u,v) \cdot (1 + \lambda \cdot |\!\sin(\theta(u,v))|)
\]
where $\theta(u,v)$ is the angle between the edge direction
$(\mathrm{centroid}(v) - \mathrm{centroid}(u))$ and $\hat{u}$.
\end{definition}

Edges running parallel to $\hat{u}$ (low $|\!\sin\theta|$) receive minimal
penalty; edges running perpendicular (high $|\!\sin\theta|$) are heavily
penalised.
This steers METIS toward cuts perpendicular to the minor axis --- straight
lines across the narrowest dimension --- and away from zigzag boundaries
parallel to the major axis.

\begin{remark}[Phase 2 deployment]
The current implementation supports $\lambda > 0$ but the empirical sweep
(Section~\ref{sec:evaluation}) uses $\lambda = 0$ (no directional penalty).
The sweep results reported here are all Phase 1 (ratio scan only).
Phase 2 is available as a CLI flag \texttt{--geosection-lambda} and its
effect on boundary straightness is deferred to future work.
\end{remark}

\subsection{METIS Implementation Details}

GeoSection calls the METIS 5.1 library via the \texttt{metis-rs} FFI crate
(vendored, no external binary required).
Key implementation choices:

\paragraph{Unequal splits: \texttt{part\_kway} vs.\ \texttt{part\_recursive}.}
For the equal-ratio case ($i = k/2$), METIS \texttt{part\_recursive} is used
(better balance for equal bisection).
For unequal ratios ($i \neq k - i$), \texttt{part\_kway} with explicit
\texttt{tpwgts = [i/k, 1-i/k]} is used.
METIS's \texttt{part\_recursive} does not support unequal vertex weight targets.

\paragraph{Per-node balance tolerance.}
At bisection node with $k$ remaining districts, the ufactor is set to
$\mathrm{ufactor} = 1 + \delta / k$, so that cumulative balance error
across all splits stays within $\delta$ per final district.
For $k=14$ (NC, first split) with $\delta=0.005$:
$\mathrm{ufactor} = 1 + 0.005/14 \approx 1.000357$ (very tight).
For $k=2$ (leaf split) with same $\delta$:
$\mathrm{ufactor} = 1 + 0.005/2 = 1.0025$ (more relaxed).

\paragraph{Post-hoc boundary-swap rebalancing.}
METIS sometimes produces $0.5$--$1$\,\% balance error without the
\texttt{-contig} option (which is incompatible with disconnected subgraphs).
A post-hoc boundary-swap loop (up to 200 iterations) moves single tracts
from the heavier side to the lighter side along the boundary until the
target balance is reached or no improvement is possible.
This resolves balance issues for all but the largest metropolitan subgraphs
(NYC, Philadelphia, Detroit).

\paragraph{Computational cost.}
For a state with $k$ districts and $N$ seeds per ratio, the top-level ratio
scan runs $\lfloor k/2 \rfloor \times N$ METIS calls.
All subsequent recursion levels use the standard single-ratio path ($N$ seeds only).
For Wisconsin ($k=8$, $N=50$): $4 \times 50 = 200$ first-level calls,
then 50 calls per node at subsequent levels.
Total METIS calls: $200 + 7 \times 50 = 550$ (the $7 = k - 1$ recursive bisections).
Compared to standard bisection ($7 \times 50 = 350$ calls), GeoSection
adds approximately $\lfloor k/2 \rfloor - 1$ ratio scans at the first level.
